% !TEX root = main.tex
\maketitle
\vspace{-1.5cm}

\section*{Approche par homogénéisation spatiale}
% ============================================================
L'approche par le module de Thiele, bien qu'utile pour une première 
analyse, repose sur des hypothèses fortes (transport purement 
diffusif, surtension uniforme, cinétique linéarisée) qui limitent 
sa portée. En particulier, elle ne permet pas de traiter rigoureusement 
le couplage entre transport et réaction dans un milieu poreux hétérogène, 
ni de faire apparaître les coefficients effectifs de manière systématique. 
L'homogénéisation spatiale permet de lever ces limitations en dérivant 
le modèle macroscopique directement depuis les équations microscopiques, 
sans hypothèse sur la forme du profil de concentration~\cite{whitakerMethodVolumeAveraging1999, leMultiscaleModelingDiffusion2017}.

La démarche est illustrée schématiquement en figure~\ref{fig:homogeneisation} : 
on part d'une description hétérogène 3D à l'échelle du pore, on applique 
la procédure de prise de moyenne pour obtenir un milieu homogène équivalent, 
puis on se ramène à un problème 1D selon la direction normale à l'électrode.

\begin{figure}[h]
\centering
\begin{tikzpicture}[
    scale=1,
    box/.style={draw, rounded corners=4pt, minimum width=2.8cm, 
                minimum height=1cm, align=center, font=\small},
    solidbox/.style={box, fill=gray!20},
    fluidbox/.style={box, fill=blue!10},
    homogbox/.style={box, fill=orange!15},
    arrow/.style={-{Stealth}, thick},
    label/.style={font=\footnotesize, align=center}
]

%--- Bloc 1 : Échelle du pore (image PDF en fond) ---
\begin{scope}[xshift=0cm]
    % Image de fond à la place du rectangle gris
    \node[inner sep=0pt, anchor=south west] at (0,0) 
        {\includegraphics[width=2.8cm,height=3.5cm]{figures/fig2_3D}};
    % Titre
    \node[below, label] at (1.4, -0.1) {3D hétérogène\\(échelle du pore $\ell_p$)};
\end{scope}

%--- Flèche + label homogénéisation ---
\begin{scope}[xshift=3.2cm]
    \draw[arrow, double, double distance=2pt] (0, 1.75) -- (1.4, 1.75);
    \node[above, font=\footnotesize, align=center] at (0.7, 1.9) 
        {Prise de};
    \node[below, font=\footnotesize, align=center] at (0.7, 1.6) 
        {moyenne};
\end{scope}


%--- Bloc 2 : VER (Volume Élémentaire Représentatif) ---
\begin{scope}[xshift=4.8cm]
    % Fond homogène
    \fill[orange!10, rounded corners=3pt] (0,0) rectangle (2.8,3.5);
    \draw[orange!60!black, thick, rounded corners=3pt] (0,0) rectangle (2.8,3.5);
    % Propriétés effectives
    \node[font=\scriptsize, align=center, orange!70!black] at (1.4, 2.6) 
        {$\varepsilon_l$, $a_v$};
    \node[font=\scriptsize, align=center] at (1.4, 2.0) 
        {$D^{eff}_i$, $\sigma^{eff}$};
    \node[font=\scriptsize, align=center] at (1.4, 1.4) 
        {$\langle C_i \rangle^l$, $\langle \varphi_l \rangle^l$};
    \node[font=\scriptsize, align=center] at (1.4, 0.8) 
        {$\langle \varphi_s \rangle^s$};
    % Titre
    \node[below, label] at (1.4, -0.1) {Milieu homogène\\équivalent};
\end{scope}

%--- Flèche + label réduction 1D ---
\begin{scope}[xshift=8.0cm]
    \draw[arrow, double, double distance=2pt] (0, 1.75) -- (1.4, 1.75);
    \node[above, font=\footnotesize, align=center] at (0.7, 1.9) 
        {Réduction};
    \node[below, font=\footnotesize, align=center] at (0.7, 1.6) 
        {1D};
\end{scope}

%--- Bloc 3 : Modèle 1D ---
\begin{scope}[xshift=9.6cm]
    % Fond 1D
    \fill[green!10, rounded corners=3pt] (0.9,0) rectangle (1.9,3.5);
    \draw[green!60!black, thick, rounded corners=3pt] (0.9,0) rectangle (1.9,3.5);
    % Axe z
    \draw[-{Stealth}, gray!60] (3,3.8) -- (3,-0.3);
    \node[right, font=\scriptsize] at (3, -0.3) {$z$};
    % Graduation en haut (paroi, z = -l_e)
    \draw[gray!60] (2.8,3.4) -- (3.2,3.4);
    \node[left, font=\tiny, gray!70!black] at (2.9,3.4) {$-\ell_e$};
    % Graduation en bas (bulk, z = 0)
    \draw[gray!60] (2.8,0) -- (3.2,0);
    \node[left, font=\tiny, gray!70!black] at (2.9,0) {$0$};
    % Conditions limites
    \node[left, font=\tiny, align=right] at (0.9, 3.4) 
        {Métal};
    \node[left, font=\tiny, align=right] at (0.9, 0.2) 
        {Phase homogène};
    % Titre
    \node[below, label] at (1.4, -0.1) {Modèle 1D\\(échelle $L$)};
\end{scope}

%--- Accolade échelles ---
\draw[decorate, decoration={brace, amplitude=6pt, mirror}, gray!60]
    (0,-1.2) -- (7.7,-1.2) 
    node[midway, below=8pt, font=\footnotesize, gray!70!black] 
    {$\ell_p \ll L$};

\end{tikzpicture}
\caption{Démarche d'homogénéisation : passage du modèle 3D hétérogène 
à l'échelle du pore au modèle 1D homogénéisé à l'échelle de l'électrode.}
\label{fig:homogeneisation}
\end{figure}





Pour cela, on applique une procédure de prise de moyenne volumique, 
qui consiste à moyenner les équations locales sur un volume élémentaire 
représentatif contenant les deux phases~\cite{whitakerMethodVolumeAveraging1999, leMultiscaleModelingDiffusion2017}. Les grandeurs locales sont 
ainsi remplacées par des grandeurs moyennées à l'échelle macroscopique, 
et les effets de la microstructure sont incorporés dans des coefficients 
effectifs. On rappelle d'abord les outils mathématiques nécessaires à 
cette démarche, avant de les appliquer successivement aux équations de 
transport et à la cinétique interfaciale.



\subsection*{Rappels théoriques}

\paragraph{Moyenne superficielle :}
On considère une grandeur $g$ définie uniquement dans le volume fluide $V_l$.
On définit sa moyenne sur le volume total $V$ (fluide + solide) par :

\begin{equation}
\avg{g} = \frac{1}{V}\int_{V_l} g\, dV \label{eq:moy_sup}
\end{equation}
Cette moyenne, dite superficielle, s’exprime comme la moyenne dans le fluide multipliée 
par la porosité $\varepsilon_l = \frac{V_l}{V}$.\cite{le<3MultiscaleModelling2022}
Elle ne correspond donc pas à la moyenne de g dans le fluide seul.
La moyenne intrinsèque $\avgl{g}$ représente la vraie valeur moyenne de $g$ 
dans le fluide seul. La relation entre moyenne superficielle et moyenne 
intrinsèque est alors :

\begin{equation}
\avg{g} = \varepsilon_l \avgl{g} = \frac{V_l}{V}\avgl{g}  \label{eq:relation_moy}
\end{equation}

\paragraph{Théorème de la divergence de Gauss :}
Ce théorème relie une intégrale de volume à une intégrale de surface. 
Pour une quantité $g$ définie dans un volume $V$ délimité par une surface $A$ :

\begin{equation}
\int_V \nab g\, dV = \oint_A \mathbf{n}\, g\, dA
\label{eq:gauss}
\end{equation}

où $\mathbf{n}$ est le vecteur normal sortant à la surface $A$.

\paragraph{Théorème de la moyenne spatiale :}
Ce théorème découle du théorème de la divergence 
de Gauss appliqué au volume fluide.
Il relie la moyenne du gradient d'une quantité $g$ à son gradient 
de moyenne, avec un terme de correction surfacique sur l'interface 
$Asf$ entre le fluide et le solide :

\begin{equation}
\avg{\nab g} = \nab\avg{ g} + \frac{1}{V}\int_{Asf} \mathbf{n}\, g\, dA \label{eq:moy_spatiale}
\end{equation}

où $\mathbf{n}$ est le vecteur normal à l'interface $Asf$ 
pointant du fluide vers le solide.

\paragraph{Théorème de transport :}
Ce théorème permet d'échanger la moyenne et la dérivée temporelle, 
à condition que le volume $V_l$ ne dépende pas du temps :

\begin{equation}
\left\langle \frac{\partial g}{\partial t} \right\rangle = \frac{\partial}{\partial t}\avg{g} \label{eq:transport}
\end{equation}

\bigskip

% ============================================================
\subsection*{Équation de conservation}

L'équation de conservation en 3D pour l'espèce $i$ s'écrit :
\begin{equation}
\frac{\partial C_i}{\partial t} + \nab \cdot N_i = 0
\label{eq:cons3D}
\end{equation}
où $C_i$ est la concentration de l'espèce $i$ et $N_i$ est le flux molaire.
En appliquant la moyenne superficielle à l'équation \eqref{eq:cons3D} :
\begin{equation}
\left\langle \frac{\partial C_i}{\partial t} \right\rangle + \avg{\nab \cdot N_i} = 0
\label{eq:cons1D}
\end{equation}
En utilisant le théorème de transport \eqref{eq:transport} et le théorème 
de la moyenne spatiale \eqref{eq:moy_spatiale} :
\begin{equation}
\frac{\partial \avg{C_i}}{\partial t} + \nab \cdot \avg{N_i} + \frac{1}{V}\int_{A_{sf}} \mathbf{n} \cdot N_i\, dA = 0
\end{equation}
Le terme surfacique $\frac{1}{V}\int_{A_{sf}} \mathbf{n}\cdot N_i\, dA$
représente l'échange de matière à l'interface fluide/solide. Comme le flux
normal est nul sur la partie non catalytique, cet échange se réduit à une
intégrale sur la seule zone active $A_{sf}^{\mathrm{cat}}$, où l'on introduit
la vitesse de réaction interfaciale $R_i$ définie par :
\begin{equation}
\mathbf{n}\cdot \mathbf{N}_i = -R_i \quad \text{sur } A_{sf}^{\mathrm{cat}}
\end{equation}
Avec la convention adoptée pour $j$, le taux $R_i$ représente la vitesse de
production algébrique de l'espèce $i$. Dans le cas d'une réaction
électrochimique, $R_i$ est relié à la densité de courant interfacial $j$ par
la loi de Faraday :
\begin{equation}
R_i =-\frac{\nu_i}{n_e F} j
\end{equation}
où $n_e$ est le nombre d’électrons échangés et $F$ la constante de Faraday.
$R_i$ porte le signe de la réaction : $R_i < 0$ si l'espèce $i$ est 
consommée à l'interface et $R_i > 0$ si elle y est produite.
 L’aire interfaciale spécifique volumique est quant à elle définie par :
\begin{equation}
a_v = \frac{A_{sf}}{V}
\end{equation} 
où $A_{sf}$ est l’aire de l’interface fluide–solide contenue dans le volume de contrôle $V$.
Cette grandeur caractérise la surface d’échange disponible par unité de volume et s’exprime en m$^{-1}$.
Le catalyseur n'étant greffé que sur une partie de cette interface, on
introduit l'aire interfaciale catalytique spécifique volumique :
\begin{equation}
a_v^{\mathrm{cat}} = \frac{A_{sf}^{\mathrm{cat}}}{V} = \theta\, a_v,
\qquad \theta = \frac{A_{sf}^{\mathrm{cat}}}{A_{sf}} \in [0,1],
\end{equation}
où $\theta$ représente la fraction catalytiquement active de l'interface.
On a donc : %Faut se mettre d'accord sur une convention pour le signe
\begin{equation}
\frac{1}{V}\int_{A_{sf}^{\mathrm{cat}}} \mathbf{n} \cdot N_i\, dA = -a_v^{\mathrm{cat}} \avg{R_i}^{sf}
\end{equation}
On obtient 
alors l'équation de conservation homogénéisée :
\begin{equation}
\varepsilon_l \frac{\partial \avgl{C_i}}{\partial t} + \nab \cdot \avg{N_i} 
= a_v^{\mathrm{cat}} \avg{R_i}^{sf}
\label{eq:cons_homo}
\end{equation}

% ============================================================
\subsection*{Flux de Nernst--Planck}

Le flux molaire de l'espèce $i$ est décrit par l'équation de Nernst-Planck, 
qui prend en compte trois contributions : diffusion, migration électrique et convection :
\begin{equation}
N_i = -D_i \nab C_i - \frac{z_i F}{RT} D_i C_i \nab \varphi_l + C_i v
\label{eq:NP3D}
\end{equation}
où $D_i$ est le coefficient de diffusion, $z_i$ la valence, $F$ la constante 
de Faraday, $R$ la constante des gaz parfaits, $T$ la température et 
$\varphi_l$ le potentiel électrique dans le fluide.

On notera que la distinction entre interface géométrique et interface
catalytique n'intervient que dans le terme source réactionnel. Pour le
transport des espèces, l'interface complète $A_{sf}$ reste imperméable,
qu'elle soit catalytique ou non.

% ============================================================
\bigskip

\subsubsection*{Terme diffusif}

En appliquant la moyenne superficielle au terme diffusif :
\begin{equation}
\avg{-D_i \nab C_i}= \frac{1}{V}\int_{V_l} (-D_i \nab C_i)\, dV \notag
\end{equation}

En appliquant le théorème de la moyenne spatiale \eqref{eq:moy_spatiale} :
\begin{equation}
\avg{-D_i \nab C_i}= -D_i \left(\nab\avg{C_i} + \frac{1}{V}\int_{A_{sf}} \mathbf{n}\, C_i\, dA \right) \notag
\end{equation}

\noindent En utilisant la relation \eqref{eq:relation_moy}, 
\noindent on a $\avg{C_i} = \varepsilon_l \avg{C_i}^l$, 
donc $\nab\avg{C_i} = \nab(\varepsilon_l \avg{C_i}^l)$. En développant :
\begin{equation}
\avg{-D_i \nab C_i}= -D_i \left(\varepsilon_l \nab\avg{C_i}^l + \avg{C_i}^l \nab\varepsilon_l
+ \frac{1}{V}\int_{A_{sf}} \mathbf{n}\, C_i\, dA \right)
\end{equation}

On décompose ensuite $C_i$ en une moyenne intrinsèque et une fluctuation à l'échelle du pore :
% ============================================================
% PAGE 2
% ============================================================
\begin{equation}
C_i = \avg{C_i}^l + \tilde{C}_i
\label{eq:decomp}
\end{equation}

où $\avg{C_i}^l$ varie à l'échelle macroscopique $l_e$ et $\tilde{C}_i$ 
varie à l'échelle du pore $\ell_p$.



\noindent En injectant la décomposition \eqref{eq:decomp} dans le terme surfacique, 
on sépare la contribution de la moyenne et celle des fluctuations :

\begin{equation}
\frac{1}{V}\int_{A_{sf}} \mathbf{n}\, C_i\, dA
= \frac{1}{V}\int_{A_{sf}} \mathbf{n}\, \avg{C_i}^l\, dA
+ \frac{1}{V}\int_{A_{sf}} \mathbf{n}\, \tilde{C}_i\, dA
\end{equation}

En appliquant le théorème de Gauss au premier terme, sachant que 
$\avg{C_i}^l$ varie à l'échelle macroscopique et peut sortir de l'intégrale :

\begin{equation}
\frac{1}{V}\int_{A_{sf}} \mathbf{n}\, \avg{C_i}^l\, dA = -\avg{C_i}^l \nab\varepsilon_l
\end{equation}

On obtient donc :
\begin{equation}
\avg{-D_i \nab C_i} = -D_i \varepsilon_l \nab\avg{C_i}^l
- D_i \frac{1}{V}\int_{A_{sf}} \mathbf{n}\, \tilde{C}_i\, dA
\end{equation}

où le premier terme correspond à la diffusion macroscopique et le second 
à la correction due aux fluctuations de concentration à l'échelle du pore.

L'équation sur $\tilde{C}_i$ est linéaire et sa source est $\nab\avg{C_i}^l$. 
On suppose que le transport à l'échelle du pore est quasi-stationnaire. La
longueur caractéristique du pore $\ell_p$ étant très petite devant la
longueur macroscopique $l_e$, le transport s'y équilibre beaucoup plus
rapidement qu'à l'échelle de l'électrode. On peut donc négliger les effets
transitoires à l'échelle du pore, ce qui donne :

\begin{equation}
\nab^2 C_i = 0 \quad \text{dans } V_l
\end{equation}

En décomposant $C_i = \avg{C_i}^l + \tilde{C}_i$ et sachant que 
$\avg{C_i}^l$ varie à l'échelle macroscopique $l_e \gg \ell_p$, 
on a $\nab^2\avg{C_i}^l \approx 0$ à l'échelle du pore. On obtient donc :

\begin{equation}
\nab^2 \tilde{C}_i = 0 \quad \text{dans } V_l
\end{equation}

avec la condition limite sur $A_{sf}$ :

\begin{equation}
-\mathbf{n} \cdot \nab\tilde{C}_i = \mathbf{n} \cdot \nab\avg{C_i}^l \quad \text{sur } A_{sf}
\end{equation}

L'équation sur $\tilde{C}_i$ est linéaire et sa source est $\nab\avg{C_i}^l$.
Par linéarité, on cherche une solution de la forme :

\begin{equation}
\tilde{C}_i = \mathbf{b}\cdot\nab\avg{C_i}^l
\end{equation}

où $\mathbf{b}$ est un vecteur de fermeture qui satisfait le problème suivant~\cite{whitakerMethodVolumeAveraging1999, korneevDataDrivenMultiscaleFramework2020} :

\begin{align}
\nab^2 \mathbf{b} &= 0 \quad \text{dans } V_l \\
-\mathbf{n}\cdot\nab\mathbf{b} &= \mathbf{n} \quad \text{sur } A_{sf}\\
\avg{\mathbf{b}}^l &= 0 \quad \text{et périodicité} : \quad b(\mathbf{r}) = b(\mathbf{r} + l_i \mathbf{e}_i)
\end{align}

En substituant cette expression dans le terme de correction surfacique, on obtient :

\begin{equation}
\frac{1}{V}\int_{A_{sf}} \mathbf{n}\, \tilde{C}_i\, dA
= \frac{1}{V}\int_{A_{sf}} \mathbf{n}\mathbf{b}\, dA\cdot\nab\avg{C_i}^l
\end{equation}

En substituant dans l'expression du terme diffusif et en factorisant :

\begin{align}
\avg{-D_i \nab C_i}
&= -D_i \left(\varepsilon_l \nab\avg{C_i}^l + \frac{1}{V}\int_{A_{sf}} \mathbf{n}\mathbf{b}\, dA\cdot \nab\avg{C_i}^l \right) \notag\\
&= -D_i \nab\avg{C_i}^l \left[\varepsilon_l I + \frac{1}{V}\int_{A_{sf}} \mathbf{n}\mathbf{b}\, dA \right] \notag\\
&= -\varepsilon_l D_i \left[ I + \frac{1}{V_l}\int_{A_{sf}} \mathbf{n}\mathbf{b}\, dA \right] \nab\avg{C_i}^l
\end{align}

On définit le coefficient de diffusion effectif $D_{\text{eff}}$ qui tient compte 
de la tortuosité du milieu poreux~\cite{whitakerMethodVolumeAveraging1999} :

\begin{equation}
D_i^{\mathrm{eff}} = D_i \left[ I + \frac{1}{V_l}\int_{A_{sf}} \mathbf{n}\mathbf{b}\, dA \right]
\end{equation}

Ce qui donne finalement :

\begin{equation}
\boxed{\avg{-D_i \nab C_i} = -\varepsilon_l D_i^{\mathrm{eff}} \nab\avg{C_i}^l}
\label{eq:diff_eff}
\end{equation}



% ============================================================
% PAGE 3
% ============================================================

% ============================================================
\subsubsection*{Terme migratoire}

Le terme migratoire fait intervenir le produit $C_i \nabla \varphi_l$, 
dont la moyenne ne se réduit pas simplement au produit des moyennes. 
On applique donc la même décomposition que pour le terme diffusif, 
en introduisant une fluctuation du potentiel électrique à l'échelle 
du pore :

\begin{equation}
\varphi_l = \avg{\varphi_l}^l + \widetilde{\varphi}_l
\end{equation}

En injectant les décompositions de $C_i$ et $\varphi_l$ dans le 
terme migratoire et en développant le produit, on fait apparaître 
quatre contributions :

\begin{align}
\left\langle -\frac{z_i F}{RT} D_i C_i \nab\varphi_l \right\rangle
&= -\frac{z_i F}{RT} D_i \avg{C_i \nab\varphi_l}
= -\frac{z_i F}{RT} D_i \frac{1}{V}\int_{V_l} C_i \nab\varphi_l\, dV \notag\\
&= -\frac{z_i F}{RT} D_i
\left[
\underbrace{\frac{1}{V}\int_{V_l} \avg{C_i}^l \nab\avg{\varphi_l}^l\, dV}_{\circled{1}}
+ \underbrace{\frac{1}{V}\int_{V_l} \avg{C_i}^l \nab\widetilde{\varphi}_l\, dV}_{\circled{2}}
\right. \notag\\
&\quad\quad\left.
+ \underbrace{\frac{1}{V}\int_{V_l} \tilde{C}_i \nab\avg{\varphi_l}^l\, dV}_{\circled{3}}
+ \underbrace{\frac{1}{V}\int_{V_l} \tilde{C}_i \nab\widetilde{\varphi}_l\, dV}_{\circled{4}}
\right]
\end{align}

On développe chacun des quatre termes :

\paragraph{Terme 1 :}
$\avg{C_i}^l$ et $\nab\avg{\varphi_l}^l$ sont constants à l'échelle du pore 
et peuvent sortir de l'intégrale :
\begin{equation}
\frac{1}{V}\int_{V_l} \avg{C_i}^l \nab\avg{\varphi_l}^l\, dV
= \avg{C_i}^l \nab\avg{\varphi_l}^l \frac{V_l}{V}
= \varepsilon_l \avg{C_i}^l \nab\avg{\varphi_l}^l
\end{equation}

\paragraph{Terme 2 :}
$\avg{C_i}^l$ est constant à l'échelle du pore et sort de l'intégrale. 
En appliquant le théorème de la moyenne spatiale \eqref{eq:moy_spatiale} 
et sachant que $\avg{\widetilde{\varphi}_l}^l = 0$ :
\begin{equation}
\frac{1}{V}\int_{V_l} \avg{C_i}^l \nab\widetilde{\varphi}_l\, dV
= \avg{C_i}^l \frac{1}{V}\int_{A_{sf}} \mathbf{n}\,\widetilde{\varphi}_l\, dA
\end{equation}

\paragraph{Terme 3 :}
$\nab\avg{\varphi_l}^l$ est constant à l'échelle du pore et sort de l'intégrale. 
Or $\avg{\tilde{C}_i}^l = 0$ par définition de la fluctuation :
\begin{equation}
\frac{1}{V}\int_{V_l} \tilde{C}_i \nab\avg{\varphi_l}^l\, dV
= \nab\avg{\varphi_l}^l \underbrace{\avg{\tilde{C}_i}^l}_{=0} = 0
\end{equation}

\paragraph{Terme 4 :}
Ce terme est un produit de deux fluctuations à l'échelle du pore. 
Par hypothèse de faible corrélation entre les fluctuations de concentration
et les fluctuations de potentiel à l'échelle du pore, ce terme est supposé
négligeable devant les contributions du premier ordre :
\begin{equation}
\frac{1}{V}\int_{V_l} \tilde{C}_i \nab\widetilde{\varphi}_l\, dV \approx 0
\end{equation}

En regroupant les termes non nuls :

\begin{equation}
\avg{C_i \nab\varphi_l}
= \varepsilon_l \avg{C_i}^l \nab\avg{\varphi_l}^l
+ \avg{C_i}^l \frac{1}{V}\int_{A_{sf}} \mathbf{n}\,\widetilde{\varphi}_l\, dA
\end{equation}

On introduit un vecteur de fermeture $\mathbf d$ permettant de
représenter les fluctuations du potentiel à l'échelle du pore sous la forme
\[
\widetilde{\varphi}_l
=
\mathbf d\cdot\nabla\langle\varphi_l\rangle^l.
\]
Cette fermeture est introduite par analogie avec le problème diffusif.
Elle constitue toutefois une approximation, puisque les fluctuations du
potentiel sont en principe couplées aux fluctuations de concentration via
l'équation de Poisson. On obtient ainsi : %MLais en fait poissson je l'ai plus là donc je le laisse mais en vrrai on sait que je peux pas garder cette hypothèse d'electroneutralité
\begin{align}
\avg{C_i \nab\varphi_l}
&= \varepsilon_l \avg{C_i}^l \nab\avg{\varphi_l}^l
+ \avg{C_i}^l \frac{1}{V}\int_{A_{sf}} \mathbf{n}\,\mathbf{d}\, dA\, \nab\avg{\varphi_l}^l \notag\\
&= \varepsilon_l \avg{C_i}^l \left(I + \frac{1}{V_l}\int_{A_{sf}} \mathbf{n}\,\mathbf{d}\, dA \right) \nab\avg{\varphi_l}^l
\end{align}

En définissant le coefficient de migration effectif $D_i^{\text{eff,mig}}$ 
de façon analogue à $D_i^{\mathrm{eff}}$ :

\begin{equation}
D_i^{\text{eff,mig}} = D_i \left(I + \frac{1}{V_l}\int_{A_{sf}} \mathbf{n}\,\mathbf{d}\, dA \right)
\end{equation}

On obtient finalement :

\begin{equation}
\boxed{-\frac{z_i F}{RT} D_i \avg{C_i \nab\varphi_l}
= -\frac{z_i F}{RT} \varepsilon_l D_i^{\text{eff,mig}} \avg{C_i}^l \nab\avg{\varphi_l}^l}
\label{eq:mig_eff}
\end{equation}

Le terme migratoire homogénéisé a ainsi la même structure que le terme 
diffusif, avec un coefficient effectif de migration $D^{\text{eff,mig}}$ 
qui tient compte de l'influence de la microstructure poreuse
(tortuosité, connectivité et géométrie locale des pores). On traite maintenant 
le terme convectif, qui introduit un phénomène supplémentaire : 
la dispersion hydrodynamique.
% ============================================================
\subsubsection*{Terme convectif}

De la même façon que pour les termes précédents, on applique la moyenne
superficielle au terme convectif. On décompose la vitesse en une moyenne
intrinsèque et une fluctuation à l'échelle du pore :
\begin{equation}
v = \avg{v}^l + \tilde{v}
\end{equation}

En injectant les décompositions de $C_i$ et $v$ :
\begin{align}
\avg{C_i v}
&= \avg{(\avg{C_i}^l + \tilde{C}_i)(\avg{v}^l + \tilde{v})} \notag\\
&= \avg{\avg{C_i}^l \avg{v}^l}
+ \avg{\avg{C_i}^l \tilde{v}}
+ \avg{\tilde{C}_i \avg{v}^l}
+ \avg{\tilde{C}_i \tilde{v}} \notag\\
&= \varepsilon_l \avg{C_i}^l \avg{v}^l
+ \avg{C_i}^l \underbrace{\avg{\tilde{v}}^l}_{=0}
+ \avg{v}^l \underbrace{\avg{\tilde{C}_i}^l}_{=0}
+ \avg{\tilde{C}_i \tilde{v}} \notag\\
&= \varepsilon_l \avg{C_i}^l \avg{v}^l
+ \avg{\tilde{C}_i \tilde{v}}.
\end{align}

Les termes $\avg{\tilde{v}}^l = 0$ et $\avg{\tilde{C}_i}^l = 0$
sont nuls par définition des fluctuations.

Le dernier terme $\avg{\tilde{C}_i \tilde{v}}$ représente la dispersion
hydrodynamique induite par les fluctuations de vitesse à l'échelle du pore.
En utilisant la fermeture obtenue précédemment,

\begin{equation}
\tilde{C}_i
=
\mathbf b \cdot \nab \avg{C_i}^l,
\end{equation}

on obtient :
\begin{align}
\avg{\tilde{C}_i \tilde{v}}
&=
\avg{
(\mathbf b \cdot \nab \avg{C_i}^l)\tilde v
}
\notag\\
&=
\avg{\mathbf b \cdot \tilde v}
\,\nab \avg{C_i}^l.
\end{align}

On définit alors le tenseur de dispersion hydrodynamique~\cite{whitakerMethodVolumeAveraging1999} :

\begin{equation}
D^{\mathrm{disp}}
=
-\avg{\mathbf b \cdot \tilde v},
\end{equation}

ce qui conduit à :

\begin{equation}
\avg{\tilde{C}_i \tilde{v}}
=
-
D^{\mathrm{disp}}
\nab \avg{C_i}^l.
\end{equation}

On obtient finalement :

\begin{equation}
\boxed{
\avg{C_i v}
=
\varepsilon_l \avg{C_i}^l \avg{v}^l
-
D^{\mathrm{disp}}
\nab \avg{C_i}^l
}
\label{eq:conv_eff}
\end{equation}

Le terme convectif homogénéisé fait ainsi apparaître une contribution
de dispersion hydrodynamique qui s'ajoute au transport convectif moyen.
Cette dispersion résulte des fluctuations de vitesse à l'échelle du pore
et se comporte comme un terme diffusif supplémentaire à l'échelle
macroscopique.

% ============================================================
% PAGE 5
% ============================================================
\subsubsection*{Réunification des termes}

En combinant les résultats obtenus pour les termes diffusif
\eqref{eq:diff_eff}, migratoire \eqref{eq:mig_eff} et convectif
\eqref{eq:conv_eff}, on obtient l'expression homogénéisée du flux de
Nernst--Planck :
\begin{align}
\avg{N_i}
&=
\varepsilon_l \avg{C_i}^l \avg{v}^l
-
\left(
\varepsilon_l D_i^{\mathrm{eff}}
+
D^{\mathrm{disp}}
\right)
\nab\avg{C_i}^l
-
\frac{z_iF}{RT}
\varepsilon_l D_i^{\mathrm{eff,mig}}
\avg{C_i}^l
\nab\avg{\varphi_l}^l .
\end{align}

On définit alors le coefficient total de transport diffusif-dispersif :
\begin{equation}
D_i^{\mathrm{tot}}
=
\varepsilon_l D_i^{\mathrm{eff}}
+
D^{\mathrm{disp}}.
\end{equation}

Le flux homogénéisé s'écrit finalement :

\begin{equation}
\boxed{
\avg{N_i}
=
\varepsilon_l \avg{C_i}^l \avg{v}^l
-
D_i^{\mathrm{tot}}
\nab\avg{C_i}^l
-
\frac{z_iF}{RT}
\varepsilon_l D_i^{\mathrm{eff,mig}}
\avg{C_i}^l
\nab\avg{\varphi_l}^l
}
\label{eq:Ni_final}
\end{equation}

% ============================================================
\subsection*{Équation de Poisson}
À titre d'extension, on présente également l'homogénéisation de 
l'équation de Poisson, qui permet de décrire explicitement les 
effets de charge d'espace lorsque l'hypothèse d'électroneutralité n'est plus vérifiée.

On réutilise ici, pour les fluctuations du potentiel $\widetilde{\varphi}_l$,
la fermeture introduite pour le terme migratoire. Cette analogie constitue
toutefois une approximation plus forte encore dans le cas de Poisson :
le problème de fermeture y est en principe couplé aux fluctuations de
concentration, puisque c'est précisément la distribution de charge qui
source le potentiel. Traiter ce problème de fermeture de façon découplée,
comme pour la diffusion ou la migration, revient à négliger ce couplage. La
dérivation qui suit doit donc être comprise comme indicative, fournissant la
structure du modèle homogénéisé plutôt qu'une fermeture rigoureuse.

En appliquant la moyenne superficielle à l'équation de Poisson :

\begin{equation}
-\varepsilon_0 \varepsilon_r \avg{\nab^2 \varphi_l} = F\sum_i z_i \avg{C_i}
\end{equation}

\subsubsection*{Membre de gauche}

En appliquant le théorème de la moyenne spatiale \eqref{eq:moy_spatiale} :

\begin{equation}
\avg{\nab\nab\varphi_l} = \nab\avg{\nab\varphi_l} + \frac{1}{V}\int_{A_{sf}} \mathbf{n}\,\nab\varphi_l\, dA
\end{equation}

La condition limite à l'interface $A_{sf}$ traduit la présence d'une 
charge de surface $\sigma_q$ :

\begin{equation}
-\varepsilon_0 \varepsilon_r \nab\varphi_l \cdot \mathbf{n} = \sigma_q 
\quad \text{sur } A_{sf}
\quad \Rightarrow \quad 
\mathbf{n}\cdot\nab\varphi_l = -\frac{\sigma_q}{\varepsilon_0 \varepsilon_r}
\end{equation}

En utilisant la fermeture introduite précédemment pour
$\widetilde{\varphi}_l$, tout en gardant à l'esprit qu'elle constitue
une approximation en présence du couplage avec l'équation de Poisson,
on obtient :
\begin{align}
\avg{\nab\nab\varphi_l}
&= \nab\avg{\nab\varphi_l} - \frac{a_v \sigma_q}{\varepsilon_0 \varepsilon_r} \notag\\
&= \nab\!\left(\nab\avg{\varphi_l}^l + \frac{1}{V}\int_{A_{sf}} \mathbf{n}\,\varphi_l\, dA\right) 
- \frac{a_v \sigma_q}{\varepsilon_0 \varepsilon_r} \notag\\
&= -\frac{a_v \sigma_q}{\varepsilon_0 \varepsilon_r}
+ \nab\!\left[\varepsilon_l \nab\avg{\varphi_l}^l
+ \frac{1}{V}\int_{A_{sf}} \mathbf{n}\,\widetilde{\varphi}_l\, dA\right] \notag\\
&= -\frac{a_v \sigma_q}{\varepsilon_0 \varepsilon_r}
+ \nab\!\left[\varepsilon_l \nab\avg{\varphi_l}^l
\left(I + \frac{1}{V_l}\int_{A_{sf}} \mathbf{n}\,\mathbf{d}\, dA\right)\right] \notag\\
&= -\frac{a_v \sigma_q}{\varepsilon_0 \varepsilon_r}
+ \nab\!\left[\varepsilon_l \frac{\varepsilon^{\text{eff}}}{\varepsilon_r} \nab\avg{\varphi_l}^l\right]
\end{align}

où $\varepsilon^{\text{eff}}$ est la permittivité effective du milieu poreux, 
définie de façon analogue à $D_{\text{eff}}$ :

\begin{equation}
\varepsilon^{\text{eff}} = \varepsilon_r \left(I + \frac{1}{V_l}\int_{A_{sf}} \mathbf{n}\,\mathbf{d}\, dA\right)
\end{equation}

\subsubsection*{Membre de droite}

En utilisant la relation entre moyenne superficielle et moyenne intrinsèque 
\eqref{eq:relation_moy} :

\begin{equation}
F\sum_i z_i \avg{C_i} = F\sum_i z_i \varepsilon_l \avg{C_i}^l
\end{equation}

\subsubsection*{Réunification}

En combinant les membres de gauche et de droite, on obtient l'équation 
de Poisson homogénéisée :

\begin{equation}
\boxed{
-\varepsilon_0 \nab\!\left(\varepsilon_l \varepsilon^{\text{eff}} \nab\avg{\varphi_l}^l\right)
= F\sum_i z_i \varepsilon_l \avg{C_i}^l + a_v \sigma_q
}
\label{eq:poisson_homo}
\end{equation}

\subsubsection*{Hypothèse d'électroneutralité}

Une alternative à l'équation de Poisson consiste à supposer 
l'électroneutralité du milieu, c'est-à-dire que la charge électrique 
nette dans le fluide est nulle :

\begin{equation}
\sum_i z_i \avg{C_i}^l = 0
\end{equation}

Cette hypothèse est justifiée lorsque la longueur de Debye reste
petite devant la taille caractéristique des pores, de sorte que les
effets de charge d'espace sont confinés au voisinage immédiat des
interfaces.
Sous cette hypothèse, la charge volumique moyenne dans l'électrolyte
est supposée négligeable. Le potentiel électrique n'est alors plus
déterminé par la résolution explicite de l'équation de Poisson, mais
par les équations de conservation du courant et les conditions aux
limites du problème.

Le potentiel dans la phase solide est quant à lui gouverné par la loi 
d'Ohm, dont la forme homogénéisée est dérivée dans la section suivante.
% ============================================================
% PAGE 6
% ============================================================
\subsection*{Loi d'Ohm}

La loi d'Ohm décrit le transport du courant électrique dans la phase
solide $\Omega_s$. Le traitement est analogue à celui du terme diffusif,
mais appliqué à la phase solide : le problème de fermeture porte
maintenant sur les fluctuations du potentiel $\widetilde{\varphi}_s$
à l'échelle du pore.


En appliquant la moyenne superficielle à l'équation de conservation du
courant dans le solide :

\begin{equation}
\avg{\nab \cdot (\sigma_s \nab\varphi_s)} = 0
\end{equation}

En notant
\begin{equation}
i_s = -\sigma_s \nab\varphi_s
\end{equation}
la densité de courant dans le solide, on obtient :

\begin{equation}
\avg{\nab \cdot i_s}=0.
\end{equation}

En appliquant le théorème de la moyenne spatiale
\eqref{eq:moy_spatiale} :

\begin{equation}
\avg{\nab \cdot i_s}
=
\nab \cdot \avg{i_s}
+
\frac{1}{V}
\int_{A_{sf}}
\mathbf n_s \cdot i_s\, dA
=
0.
\end{equation}

En identifiant le terme surfacique au courant échangé à l'interface
fluide/solide, on obtient :
\begin{equation}
\nab \cdot \avg{i_s}
=
-a_v^{\mathrm{cat}} \avg{j}^{sf}.
\end{equation}

où $\avg{j}^{sf}$ est la densité de courant échangée à l'interface
fluide/solide.

On cherche maintenant à relier $\avg{i_s}$ au potentiel moyen dans la
phase solide. En développant
$\avg{i_s} = -\sigma_s \avg{\nab\varphi_s}$ et en appliquant le théorème
de la moyenne spatiale :

\begin{align}
\avg{\nab\varphi_s}
&= \nab\avg{\varphi_s} + \frac{1}{V}\int_{A_{sf}} \mathbf{n}_s\, \varphi_s\, dA \notag\\
&= \varepsilon_s \nab\avg{\varphi_s}^s + \frac{1}{V}\int_{A_{sf}} \mathbf{n}_s\, \varphi_s\, dA
\end{align}

Comme pour le terme diffusif, les fluctuations du potentiel solide sont
représentées à l'aide d'un vecteur de fermeture $e$ défini par

\begin{equation}
\widetilde{\varphi}_s
=
e\cdot\nab\avg{\varphi_s}^s.
\end{equation}

Le vecteur de fermeture $e$ satisfait alors le problème local :
\begin{align}
\nab^2 e &= 0 && \text{dans } V_s,\\
-\mathbf n_s\cdot\nab e &= \mathbf n_s && \text{sur } A_{sf},\\
\avg e^s &=0,
\end{align}
complété par des conditions de périodicité sur les frontières de la cellule élémentaire.

En substituant dans l'expression de $\avg{\nab\varphi_s}$ :

\begin{equation}
\frac{1}{V}\int_{A_{sf}} \mathbf{n}_s \widetilde{\varphi}_s\, dA
= \varepsilon_s \frac{1}{V_s} \int_{A_{sf}} \mathbf{n}_s\, e\, dA\, \nab\avg{\varphi_s}^s
\end{equation}

On définit la conductivité effective $\sigma^{\text{eff}}$ du milieu poreux :

\begin{equation}
\sigma^{\text{eff}} = \sigma_s \left(I + \frac{1}{V_s}\int_{A_{sf}} \mathbf{n}_s\, e\, dA \right)
\end{equation}

D'où :

\begin{equation}
\avg{\nab\varphi_s}
= \varepsilon_s \left(I + \frac{1}{V_s}\int_{A_{sf}} \mathbf{n}_s\, e\, dA \right) \nab\avg{\varphi_s}^s
= \varepsilon_s \frac{\sigma^{\text{eff}}}{\sigma_s} \nab\avg{\varphi_s}^s
\end{equation}

On obtient finalement la loi d'Ohm homogénéisée :

\begin{equation}
\boxed{\nab\!\left(\varepsilon_s \sigma^{\text{eff}} \nab\avg{\varphi_s}^s\right) = a_v^{\mathrm{cat}} \avg{j}^{sf}}
\label{eq:ohm_homo}
\end{equation}

L'homogénéisation des équations de transport, de l'équation de Poisson
et de la loi d'Ohm conduit ainsi à un système macroscopique décrivant
les transferts de matière et de charge dans l'électrode poreuse~\cite{mendez-vasquezUpscalingMassElectric2025}. Les
effets de la microstructure sont incorporés dans les différents
coefficients effectifs, ce qui permet de représenter leur influence sans
résoudre explicitement la géométrie des pores.

Ce système constitue le cadre général de la modélisation 1D. Différentes
hypothèses simplificatrices peuvent ensuite être introduites selon les
échelles caractéristiques du problème et les phénomènes physiques que
l'on souhaite décrire. Les conséquences de ces hypothèses sur les
équations gouvernantes seront examinées dans la section suivante.

